\chapter{Fractional Sobolev Norm}

We will now explore a different method to price American puts with neural networks, and this time we will work with Fractional Sobolev Norms to construct our loss function. The benefit of this is that we would be able to guarantee that our neural network converges to the true solution provided that its loss converges, and that our loss conditions apply for all situations and is not dependent on the payoff.

\section{Preliminaries}

\subsection{Weak derivatives}

Suppose $u \in L_\text{loc}^1(\Omega)$, a locally integrable function on $\Omega$. This means that $u \in L^1(U)$ for every open set $U$ such that the closure $\overline U \subset \Omega$ and $\overline U$ is a compact (closed and bounded) subset of $\mathbb R^n$.

[finish later]

\subsection{Sobolev spaces}

Here, we introduce Sobolev spaces of integer order, defined over an arbitrary domain $\Omega \subset \mathbb R^n$. These are vector subspaces of various Lebesgue spaces $L^p(\Omega)$, which is the class of all measurable functions $u$ for which
$$\int_\Omega |u(x)|^p \,\mathrm dx < \infty$$

We first define the Sobolev norms, which is a functional $\|\cdot\|_{m, p}$ where $m \in \mathbb N_+$ and $1 \le p \le \infty$, as follows:
\begin{align*}
    \|u\|_{m, p} &= \left(\sum_{0 \le |\alpha|\le m} \|D^\alpha u\|^p_p\right)^{1/p} \qquad  1 \le p < \infty \\
    \|u\|_{m, \infty} &= \max_{0 \le |\alpha| \le m }\|D^\alpha u\|_\infty
\end{align*}

where $\|\cdot\|_p$ is the norm in $L^p(\Omega)$. The Sobolev spaces are consequently those on which $\|\cdot \|_{m, p}$ is a norm \cite{adams2003sobolev}:
\begin{itemize}
    \item $H^{m, p}(\Omega) \equiv \text{the completion of } \{u \in C^m (\Omega) : \|u\|_{m, p} < \infty\}$ with respect to the norm $\|\cdot \|_{m, p}$
    \item $W^{m, p}(\Omega) \equiv \{u \in L^p(\Omega) : D^\alpha u \in L^p(\Omega) \text{ for } 0 \le |\alpha| \le m \}$ where $D^\alpha u$ is the weak partial derivative
    \item $W_0^{m, p}(\Omega) \equiv \text{the closure of }C^\infty_0 \text{ in the space } W^{m, p}(\Omega)$
\end{itemize}

It can be shown by Theorem 3.17 of \cite{adams2003sobolev} that for $1 \le p < \infty$ we have
$$H^{m, p}(\Omega) = W^{m, p}(\Omega)$$
so for the rest of this section, we will use the space $H^{m, p}(\Omega)$. 

\subsection{Trace theory}




\subsection{Fractional Sobolev spaces}
Suppose we have a general, possibly nonsmooth, open set $\Omega \subset \mathbb R^n$. For any real $s > 0$ and for any $p \in [1, \infty)$, we aim to define the fractional Sobolev spaces $W^{s, p}(\Omega)$.

To construct this, let's first fix the fractional exponent $s \in (0, 1)$. Then for any $p \in [1, +\infty)$ we define $W^{s, p}(\Omega)$ as \cite{di2012hitchhikers}:
$$W^{s, p}(\Omega) := \left\{u \in L^p(\Omega) : \frac{|u(x) - u(y)|}{|x-y|^{\frac np + s}} \in L^p(\Omega \times \Omega)\right\}$$

which can be interpreted as an intermediary Banach space between $L^p(\Omega)$ and $W^{1, p}(\Omega)$ endowed with the norm

$$\|u\|_{W^{s, p}(\Omega)} := \left(\int_\Omega |u|^p \,\mathrm dx + \int_\Omega \int_\Omega \frac{|u(x)-u(y)|^p}{|x-y|^{n+sp}}\,\mathrm dx \mathrm dy\right)^{\frac 1p}$$
where we have used the \textit{Gagliardo seminorm} of $u$:
$$[u]_{W^{s, p}(\Omega)} := \left( \int_\Omega \int_\Omega \frac{|u(x)-u(y)|^p}{|x-y|^{n+sp}}\,\mathrm dx \mathrm dy\right)^{\frac 1p}$$


\section{Model}
Suppose we would like to price an option on $k$ assets $\bm S = (S^1, \dots, S^k)$. Similar to before, these assets have theoretically unbounded prices, so we need to truncate it a minimum and maximum price in our model. Thus the domain of the target function can be chosen to be

$$\Omega = (S_\text{min}^1, S_\text{max}^1) \times \cdots \times (S_\text{min}^k, S_\text{max}^k)$$

In our study we will make the simplification that $S_\text{min}^i = S_\text{min}^j$ and $S_\text{max}^i = S_\text{max}^j$ for all $i \ne j$, which we can call $a$ and $b$ respectively. Our spaces would consequently have the definitions

$$\Omega = (a, b)^n \qquad \Omega_T = (0, T) \times \Omega \qquad \Sigma = (0, T) \times \partial \Omega$$

where $\partial \Omega =: \Gamma$ is the boundary of $\Omega$, a hyper-rectangle. We will use the representation

$$\Gamma = \bigcup_{i=1}^n \big(\{x \in [a, b]^n : x_i = a\} \cup \{x \in [a, b]^n : x_i = b\} \big)$$
which is not the simplest, but it is useful as it is decomposed into faces, which is useful when we would like to use the norm derivative later. Note that in the one-dimensional case, this reduces simply to the points $\{a, b\}$.

\section{Loss function}
Again, our loss function would be a weighted sum of four loss components, each of which ensures regulation of a specific aspect of our function.

\begin{enumerate}
    \item The variational loss, which is as before, controlling the accuracy of our function in the interior of our domain:
    $$J_1 = \|\min \{\mathcal G[f], f-g\}\|^2_{L^2(\Omega_T), \mu_1}$$
    \item The energy loss: before we were only controlling the L2 norm of the payoff, but we require the option value to equal the payoff as a function and not just pointwise. Thus we need to incorporate the gradient into the loss so that the spatial slope of the payoff is respected.
    $$J_2 = \|f(0, \cdot) - g(\cdot) \|^2_{H^1(\Omega), \mu_2}$$
    where we have the norm defined as
    $$\|d\|^2_{H^1(\Omega), \mu_2} = \|d\|^2_{L^2(\Omega)} + \|\nabla d\|^2_{L^2(\Omega)}$$
    and the derivative
    $$\nabla d = \big(\partial_{x_1}d, \dots, \partial_{x_k}d\big)$$
    \item J3 (need to explain intuitively)
    $$J_3 =\|f(t, x) - g( x)\|^2_{H^{\frac 34, \frac 32}(\Sigma), \mu_3}$$
    where
    \begin{align*}
        \|d\|^2_{H^{\frac 34, \frac 32}(\Sigma)} = \|d\|^2_{H^{0,1}(\Sigma)} &+ \int_0^T\int_0^T\frac{\|d(t_1, \cdot) - d(t_2, \cdot)\|^2_{L^2(\Gamma)}}{|t_1-t_2|^{2.5}}\,\mathrm dt_1\,\mathrm dt_2 \\ &+ \int_0^T\int_\Gamma\int_\Gamma \frac{|d(t, x) - d(t, y)|^2}{\|x-y\|^{2}}\,\mathrm d\sigma(x)\mathrm d\sigma(y) \mathrm dt
    \end{align*}
    The time exponent $2.5$ is equal to $2p+1$ where $p$ is the decimal part of the time fractional exponent, and the space exponent $2$ is equal to $2\theta + k - 1$ where $\theta$ is the decimal part of the space fractional exponent, and $k$ the dimension. Note that here $\mathrm d \sigma(x)$ and $\mathrm d \sigma(y)$ indicate the surface measures on $\Gamma$, but in practice we can treat these as $\mathrm dx$ and $\mathrm dy$.
    \item J4 (need to explain intuitively):
    $$J_4 =\left\|\nabla (f-g) \cdot v\right\|^2_{H^{\frac 14, \frac 12}(\Sigma), \mu_4}$$
    where the norm derivative is $\nabla(f-g)\cdot \nu = \partial_{x_{i^\star}}(f-g)$ given that $x_{i^\star} \in \{a, b\}$ and $x_{i} \in (a, b)$ for $i \ne i^\star$. Rigorously, we require the sign but in practice this is not necessary, as we will square in the end. The fractional norm is thus written
    \begin{align*}
        \|d\|^2_{H^{\frac 14, \frac 12}(\Sigma)} = \|d\|^2_{L^2(\Sigma)} &+ \int_0^T\int_0^T \frac{\|d(t_1, \cdot) - d(t_2, \cdot)\|^2_{L^2(\Gamma)}}{|t_1-t_2|^{1.5}}\,\mathrm dt_1\mathrm dt_2 \\ &+ \int_{0}^T\int_\Gamma\int_\Gamma \frac{|d(t, x) - d(t, y)|^2}{\|x-y\|^2}\,\mathrm d\sigma(x)\mathrm d(y)\mathrm dt
    \end{align*}
    because $H^{0, 0} = L^2$.
    
\end{enumerate}

In practice, we will compute a Monte Carlo estimate of these integrals by sampling pairs $(t_1, t_2)$ from the domain, evaluating the integrand and taking the mean across all pairs. To see this, suppose we can sample iid a random variable $X$ with probability density function $q(x)$. Then for any integrand $\phi(x)$ we have
    
$$
\int \phi(x)\,\mathrm dx = \int\frac{\phi(x)}{q(x)}q(x) \,\mathrm dx = \mathbb E\left[\frac{\phi(X)}{q(X)}\right] \approx \frac 1N\sum_{i=1}^N\frac{\phi(X_i)}{q(X_i)}
$$

which can be generalised similarly to higher dimensions. By restricting our domain of $t_1$ and $t_2$ to be uniform in $[0, 1]$, we have $q$ uniformly equal to $1$ and hence we can just evaluate the integrand.


\section{Numerical results}
